\rok{Јануар 2026.}{А}

Prvi poprаvni test iz Algoritama i strukture podataka

Prvi praktični test održan 29.01.2026.

\zadatak{1}{Selection Sort}
Написати Selection Sort алгоритам за сортирање целобројног низа дужине n.

\textbf{Додатне напомене за решавање задатка:}
\begin{itemize}
	\item Улаз је несортиран низ целих бројева
	\item Излаз је сортиран низ целих бројева
	\item У template-у се налази класа \texttt{RandomArrayGenerator} коју треба искратити за генерисање низа случајних бројева
	\item Креирати класу Selection и у оквиру ње генерисати методу:
	\begin{Verbatim}[fontsize=\small]
public static void Sort(long[] arr)
	\end{Verbatim}
\end{itemize}

\zadatak{2}{Сума елемената између првог и последњег појављивања максималне вредности}
Имплементирати методу која проналази максималну вредност у листи. Након тога, потребно је израчунати и вратити суму свих елемената који се налазе између првог и последњег појављивања те максималне вредности (укључујући и те граничне елементе).

\textbf{Додатни захтеви за решавање задатка:}
\begin{itemize}
	\item Ако се максимална вредност појављује само једном, сума је једнака тој самој вредности.
	\item Алгоритам писати у оквиру класе LinkedList.
	\item Ако листа не постоји, вратити 0.
	\item Алгоритам мора имати временску комплексност O(n).
	\item Захтеван потпис методе:
	\begin{Verbatim}[fontsize=\small]
public int sumBetweenMaxBounds()
	\end{Verbatim}
\end{itemize}

\textbf{Примери:}
\begin{itemize}
	\item \textbf{Улаз 1:} \texttt{4 → 2 → 3 → 4}
	\item \textbf{Излаз 1:} \texttt{13}
	\item \textbf{Улаз 2:} \texttt{2 → 3 → 4 → 4}
	\item \textbf{Излаз 2:} \texttt{8}
	\item \textbf{Улаз 3:} \texttt{2 → 3 → 1 → 4}
	\item \textbf{Излаз 3:} \texttt{0}
\end{itemize}

\zadatak{3}{Гравитациона Дробилица (The Gravity Crusher)}
Замисли да имаш вертикалну cev у коју убацујеш блокове различитих тежина. Сваки блок има своју тежину (позитиван цео број). Када нови блок падне на онај испод њега, важе следећа правила:

\textbf{Правила:}
\begin{itemize}
	\item \textbf{Дробљење:} Ако је нови блок строго тежи од блока на врху, он ће га смрскати (уклонити из стека) и покушати да падне на следећи блок испод њега.
	\item \textbf{Пуцање:} Ако је нови блок исте тежине као блок на врху, оба блока пуцају и оба нестају.
	\item \textbf{Стабилност:} Ако је нови блок лакши од блока на врху, он једноставно остаје да лежи на њему.
\end{itemize}

\textbf{Твој циљ:} Имплементирати методу која прима низ тежина које се редом убацују у cev и враћа финално стање блокова у cevi (од дна ка врху).

\textbf{Додатни захтеви за решавање задатка:}
\begin{itemize}
	\item Временска сложеност: O(n).
	\item Резултат мора бити низ у оригималном редоследу (од блока који је убачен и остао на дну, до последњег на врху).
\end{itemize}

\textbf{Примери:}
\begin{itemize}
	\item \textbf{Улаз 1:} \texttt{[3, 10, 2, 2, 8]}
	\item \textbf{Излаз 1:} \texttt{[10, 8]}
	\item \textbf{Улаз 2:} \texttt{[7, 6, 10, 5, 4]}
	\item \textbf{Излаз 2:} \texttt{[10, 5, 4]}
\end{itemize}

\sablon{Шаблон првог поправног теста јануар 2026. група А}{2025/Januar/GrupaA/AISP2025_PopravniTest1_Test1_GrupaA.linq}
